\documentclass[final]{beamer}

% ------------------------------------------------------------
% Packages
% ------------------------------------------------------------

\usepackage[size=a1,orientation=portrait,scale=1.08]{beamerposter}
\usepackage{amsmath,amssymb,mathtools}
\usepackage{tikz}
\usepackage{tikz-cd}
\usepackage{booktabs}
\usepackage{ragged2e}
\usepackage{microtype}
\usepackage{xcolor}

\usepackage{mlmodern}
% ------------------------------------------------------------
% Colours
% ------------------------------------------------------------

\definecolor{PosterBlue}{RGB}{28,65,110}
\definecolor{PosterLightBlue}{RGB}{232,240,248}
\definecolor{PosterDark}{RGB}{35,35,35}
\definecolor{PosterAccent}{RGB}{166,69,54}
\definecolor{PosterGrey}{RGB}{245,245,245}

% ------------------------------------------------------------
% Beamer style
% ------------------------------------------------------------

\setbeamercolor{normal text}{
  fg=PosterDark,
  bg=white
}

\setbeamercolor{headline}{
  fg=white,
  bg=PosterBlue
}

\setbeamercolor{block title}{
  fg=white,
  bg=PosterBlue
}

\setbeamercolor{block body}{
  fg=PosterDark,
  bg=PosterGrey
}

\setbeamercolor{block title alerted}{
  fg=white,
  bg=PosterAccent
}

\setbeamercolor{block body alerted}{
  fg=PosterDark,
  bg=PosterLightBlue
}

\setbeamercolor{itemize item}{
  fg=PosterAccent
}

\setbeamercolor{itemize subitem}{
  fg=PosterBlue
}

\setbeamertemplate{navigation symbols}{}
\setbeamertemplate{itemize item}{\raisebox{0.15em}{\small$\bullet$}}
\setbeamertemplate{itemize subitem}{\raisebox{0.1em}{\scriptsize$\bullet$}}

\setbeamertemplate{block begin}{
  \vskip1ex
  \begin{beamercolorbox}[
    rounded=true,
    shadow=false,
    sep=1.2ex,
    wd=\linewidth
  ]{block title}
    \usebeamerfont{block title}\insertblocktitle
  \end{beamercolorbox}
  \nointerlineskip
  \begin{beamercolorbox}[
    rounded=true,
    shadow=false,
    sep=1.5ex,
    wd=\linewidth
  ]{block body}
}

\setbeamertemplate{block end}{
  \end{beamercolorbox}
  \vskip1ex
}

\setbeamertemplate{block alerted begin}{
  \vskip1ex
  \begin{beamercolorbox}[
    rounded=true,
    shadow=false,
    sep=1.2ex,
    wd=\linewidth
  ]{block title alerted}
    \usebeamerfont{block title}\insertblocktitle
  \end{beamercolorbox}
  \nointerlineskip
  \begin{beamercolorbox}[
    rounded=true,
    shadow=false,
    sep=1.5ex,
    wd=\linewidth
  ]{block body alerted}
}

\setbeamertemplate{block alerted end}{
  \end{beamercolorbox}
  \vskip1ex
}

% ------------------------------------------------------------
% Typography
% ------------------------------------------------------------

\setbeamerfont{headline title}{
  size=\veryHuge,
  series=\bfseries
}

\setbeamerfont{headline author}{
  size=\Large
}

\setbeamerfont{headline institute}{
  size=\large
}

\setbeamerfont{block title}{
  size=\LARGE,
  series=\bfseries
}

\setbeamerfont{block title alerted}{
  size=\LARGE,
  series=\bfseries
}

\setbeamerfont{normal text}{
  size=\large
}

% ------------------------------------------------------------
% Shortcuts
% ------------------------------------------------------------

\newcommand{\DgmSet}{\mathbf{DgmSet}}
\newcommand{\RDC}{\mathbf{RDC}}
\newcommand{\CatInf}[1]{(\infty,#1)\text{-}\mathbf{Cat}}
\newcommand{\Gray}{\otimes}
\newcommand{\op}{\mathrm{op}}

\newcommand{\placeholder}[1]{%
  \textcolor{PosterAccent}{\textbf{[#1]}}%
}

% ------------------------------------------------------------
% Poster metadata
% ------------------------------------------------------------

\title{The Diagrammatic Model of \texorpdfstring{$(\infty,n)$}{(infinity,n)}-Categories}

\author{
  Clémence Chanavat
  \qquad
  \textnormal{joint work with Amar Hadzihasanovic}
}

\institute{
  Tallinn University of Technology
  \qquad
  \texttt{\placeholder{email address}}
}

% ------------------------------------------------------------
% Document
% ------------------------------------------------------------

\begin{document}

\begin{frame}[fragile, t]

% ------------------------------------------------------------
% Header
% ------------------------------------------------------------

\begin{beamercolorbox}[
  wd=\paperwidth,
  sep=2.2ex,
  center
]{headline}

  {\usebeamerfont{headline title}\inserttitle\par}

  \vspace{0.6ex}

  {\usebeamerfont{headline author}\insertauthor\par}

  \vspace{0.3ex}

  {\usebeamerfont{headline institute}\insertinstitute\par}

\end{beamercolorbox}

\vspace{1.5ex}

% ------------------------------------------------------------
% Columns
% ------------------------------------------------------------

\begin{columns}[t,totalwidth=\textwidth]

% ============================================================
% COLUMN 1
% ============================================================

\begin{column}{0.315\textwidth}

\begin{block}{Higher categories and strictification}

Higher categories encode structures whose composition laws are associative
and unital only up to coherent higher equivalences.

\medskip

Familiar examples include:

\begin{itemize}
  \item bicategories and tricategories;
  \item quasicategories and complete Segal spaces;
  \item higher categorical structures arising in topology,
        algebra, and mathematical physics.
\end{itemize}

\medskip

A fundamental question is:

\begin{center}
\Large
\textbf{How much of a homotopy-coherent higher category can be made strict?}
\end{center}

\medskip

In general, an $(\infty,n)$-category cannot be replaced by a strict
$n$-category. The diagrammatic model instead admits an intermediate
form of strictification in which composition remains algebraic while
exchange and associativity become strict.

\end{block}

\begin{block}{Diagrammatic shapes}

The model is based on a category of finite shapes constructed from
\emph{regular directed complexes}.

\medskip

Informally, these are regular CW-complexes equipped with compatible
input and output boundaries in every dimension.

\medskip

A diagrammatic set is a presheaf
\[
  X \colon \mathcal{A}^{\op}\longrightarrow \mathbf{Set},
\]
where $\mathcal{A}$ is a category of atomic directed shapes.

\medskip

For an atom $U$, an element
\[
  x\in X(U)
\]
is interpreted as a diagram of shape $U$ in $X$.

\medskip

\begin{center}
\begin{tikzpicture}[
  scale=1.4,
  every node/.style={font=\large},
  vertex/.style={
    circle,
    fill=PosterBlue,
    inner sep=3pt
  },
  arr/.style={
    ->,
    very thick,
    PosterDark
  }
]
  \node[vertex] (a) at (0,0) {};
  \node[vertex] (b) at (2,0) {};
  \node[vertex] (c) at (4,0) {};

  \draw[arr,bend left=25] (a) to node[above] {$f$} (b);
  \draw[arr,bend left=25] (b) to node[above] {$g$} (c);
  \draw[arr,bend right=28] (a) to node[below] {$h$} (c);

  \draw[
    ->,
    very thick,
    PosterAccent
  ]
  (2,0.55) -- (2,1.6)
  node[midway,right] {$\alpha$};
\end{tikzpicture}
\end{center}

The available shapes are not restricted to globes or simplices:
they include arbitrary finite pasting diagrams satisfying suitable
regularity conditions.

\end{block}

\begin{block}{Why use arbitrary pasting diagrams?}

The geometry of the shapes directly encodes composition.

\begin{itemize}
  \item Composite diagrams are represented by actual glued cell complexes.
  \item Higher-dimensional pasting is native rather than reconstructed indirectly.
  \item Diagrammatic operations can be defined directly on shapes.
  \item Algebraic composition can coexist with homotopy-coherent equivalence data.
\end{itemize}

\end{block}

\end{column}

% ============================================================
% COLUMN 2
% ============================================================

\begin{column}{0.35\textwidth}

\begin{block}{Native operations}

The category of diagrammatic sets supports a wide collection of
higher-categorical constructions.

\begin{center}
\renewcommand{\arraystretch}{1.4}
\begin{tabular}{@{}ll@{}}
\toprule
\textbf{Operation} & \textbf{Geometric meaning} \\
\midrule
Gray product $X\Gray Y$
  & directed interaction of diagrams \\

Join $X\star Y$
  & adjoining diagrams in successive dimensions \\

Suspension $\Sigma X$
  & shifting the dimension of cells \\

Opposites $X^{\op}$
  & reversing selected directions \\

Markings
  & recording equivalences algebraically \\
\bottomrule
\end{tabular}
\end{center}

\medskip

These constructions are defined from corresponding operations on the
underlying directed cell complexes.

\medskip

In particular, the Gray product reflects genuinely directed
higher-dimensional interaction, rather than a cartesian product of
underlying sets.

\end{block}

\begin{block}{Equivalences}

The model admits both inductive and coinductive notions of equivalence.

\medskip

Very roughly, a cell is an equivalence if it admits inverse data,
together with higher cells witnessing the required equations.

\[
\begin{tikzcd}[column sep=huge,row sep=large]
x
  \arrow[r,bend left=18,"f"]
&
y
  \arrow[l,bend left=18,"g"]
\end{tikzcd}
\]

together with higher cells of the form
\[
  g\circ f \Rightarrow \mathrm{id}_x,
  \qquad
  f\circ g \Rightarrow \mathrm{id}_y,
\]
and further coherence data.

\medskip

The coinductive definition allows these witnesses themselves to be
equivalences, recursively in all dimensions.

\end{block}

\begin{alertblock}{Semi-strictification theorem}

\begin{center}
\Large
\textbf{
Every diagrammatic $(\infty,n)$-category is equivalent to one whose
composition laws are algebraic and satisfy strict associativity and
strict exchange.
}
\end{center}

\medskip

Thus the model separates two sources of weakness:

\begin{itemize}
  \item \textbf{composition}, which can be made algebraic and strictly compatible;
  \item \textbf{equivalence data}, which remains genuinely higher-dimensional
        and homotopy-coherent.
\end{itemize}

\medskip

Schematically,
\[
  X
  \simeq
  X_{\mathrm{ss}},
\]
where $X_{\mathrm{ss}}$ is semi-strict.

\medskip

This is stronger than merely choosing composites, but weaker than
replacing $X$ by a strict $n$-category.

\end{alertblock}

\begin{block}{What becomes strict?}

For composable diagrams $u,v,w$, one obtains equations such as
\[
  (u\circ v)\circ w
  =
  u\circ(v\circ w)
\]
on the nose.

For independent compositions, exchange is also strict:
\[
  (\alpha\circ_1\beta)\circ_0(\gamma\circ_1\delta)
  =
  (\alpha\circ_0\gamma)\circ_1(\beta\circ_0\delta).
\]

\medskip

The theorem does \emph{not} assert that all higher cells are uniquely
determined, or that all equivalences become identities.

\end{block}

\end{column}

% ============================================================
% COLUMN 3
% ============================================================

\begin{column}{0.315\textwidth}

\begin{block}{Model structures}

For every $n\in\mathbb{N}\cup\{\infty\}$, diagrammatic sets admit a
model structure presenting $(\infty,n)$-categories.

\medskip

The fibrant objects satisfy horn-filling and equivalence-completeness
conditions adapted to arbitrary directed pasting shapes.

\medskip

The weak equivalences encode the expected homotopy theory of
$(\infty,n)$-categories.

\medskip

There are also marked versions in which equivalences are specified
as part of the structure.

\begin{center}
\[
\begin{tikzcd}[column sep=large]
\text{marked diagrammatic sets}
  \arrow[r,shift left=1.2ex]
&
\text{diagrammatic sets}
  \arrow[l,shift left=1.2ex]
\end{tikzcd}
\]
\end{center}

\end{block}

\begin{block}{The Gray product}

A key structural result is that the relevant model structures are
compatible with the Gray tensor product.

\medskip

This permits the construction of:

\begin{itemize}
  \item enriched higher categories;
  \item internal hom objects;
  \item cylinder and path constructions;
  \item monoidal homotopy-theoretic operations.
\end{itemize}

\medskip

At the level of shapes, the Gray product combines directed cells while
remembering the order in which their directions interact.

\begin{center}
\begin{tikzpicture}[
  scale=1.25,
  every node/.style={font=\large},
  arr/.style={->,very thick}
]
  \draw[arr] (0,0) -- (2.3,0);
  \draw[arr] (0,0) -- (0,2.3);
  \draw[arr] (0,2.3) -- (2.3,2.3);
  \draw[arr] (2.3,0) -- (2.3,2.3);

  \fill[PosterBlue] (0,0) circle (3pt);
  \fill[PosterBlue] (2.3,0) circle (3pt);
  \fill[PosterBlue] (0,2.3) circle (3pt);
  \fill[PosterBlue] (2.3,2.3) circle (3pt);

  \draw[
    very thick,
    PosterAccent,
    ->,
    bend left=18
  ]
  (0.45,0.35) to (1.85,1.95);

  \node at (1.15,-0.35) {$f$};
  \node[rotate=90] at (-0.35,1.15) {$g$};
  \node at (1.2,1.15) {$f\Gray g$};
\end{tikzpicture}
\end{center}

\end{block}

\begin{block}{Outlook}

The diagrammatic model provides a setting in which one can study:

\begin{itemize}
  \item comparisons with simplicial and stratified models;
  \item strictification and pasting theorems;
  \item general-position methods for directed cell complexes;
  \item factorization homology and stratified topology;
  \item computational representations of higher categorical diagrams.
\end{itemize}

\medskip

A major open direction is to understand more explicitly how
diagrammatic shapes compare with the shapes underlying other models
of higher categories.

\end{block}

\begin{block}{References}

\small

\begin{thebibliography}{9}

\bibitem{hadzihasanovic}
A.~Hadzihasanovic,
\newblock \emph{\placeholder{Title of foundational diagrammatic-sets paper}},
\newblock \placeholder{journal or arXiv reference}.

\bibitem{equivalences}
C.~Chanavat and A.~Hadzihasanovic,
\newblock \emph{Equivalences in diagrammatic sets},
\newblock Journal of Pure and Applied Algebra,
230(1), 2026.

\bibitem{pasting}
C.~Chanavat,
\newblock \emph{\placeholder{Title of the semi-strictification or pasting paper}},
\newblock \placeholder{arXiv reference}.

\bibitem{gray}
C.~Chanavat and A.~Hadzihasanovic,
\newblock \emph{\placeholder{Title of the Gray-product paper}},
\newblock \placeholder{arXiv reference}.

\end{thebibliography}

\end{block}

\begin{block}{Contact}

\begin{columns}[c,totalwidth=\linewidth]

\begin{column}{0.68\linewidth}
\large

\textbf{Clémence Chanavat}

\medskip

Tallinn University of Technology

\medskip

\texttt{\placeholder{email}}

\medskip

\texttt{\placeholder{website or arXiv profile}}

\end{column}

\begin{column}{0.27\linewidth}
\centering

\fbox{
  \begin{minipage}[c][7cm][c]{7cm}
    \centering
    \textbf{QR code}

    \medskip

    \placeholder{insert QR code}
  \end{minipage}
}

\end{column}

\end{columns}

\end{block}

\end{column}

\end{columns}

\end{frame}

\end{document}
